
\chapter{Characterizing extremely amenable automorphism groups by a Ramsey property}


\section{Dynamics and Extreme Amenability}
\label{sec:dynamics_extreme_amenability}

We will first reformulate the condition that a $G$-flow has a fixed point
in the following manner.


\begin{lemma}
    \label{lem:fixed_point_equiv}
    Let $G$ be a topological group and let $X$ be a $G$-flow. Then the following are equivalent:
    \begin{enumerate}
        \item[(i)] The $G$-flow $X$ has a fixed point.
        \item[(ii)] For every integer $n \ge 1$, every continuous function $f: X \to \mathbb{R}^n$, every $\epsilon > 0$, and every finite subset $F \subseteq G$, there exists $x \in X$ such that $\|f(x) - f(g \cdot x)\| \le \epsilon$ for all $g \in F$, where $\|\cdot\|$ denotes the Euclidean norm.
    \end{enumerate}
\end{lemma}

\begin{proof}
    (i) $\implies$ (ii). This is immediate from the definition of a fixed point.

    (ii) $\implies$ (i). We proceed via a compactness argument. For any continuous function $f: X \to \mathbb{R}^n$, $\epsilon > 0$, and finite subset $F \subseteq G$, define the set
    $$
        A_{f, \epsilon, F} = \{x \in X : \|f(x) - f(g \cdot x)\| \le \epsilon \text{ for all } g \in F\}.
    $$

    \textbf{Claim.} The intersection $\bigcap_{f, \epsilon, F} A_{f, \epsilon, F}$ is non-empty.

    \textit{Proof of the Claim.} Because $f$ and the group action are continuous, each set $A_{f, \epsilon, F}$ is closed. Since $X$ is compact, it suffices to show that this collection of closed sets has the finite intersection property. Let $(f_j, \epsilon_j, F_j)$ for $j = 1, \dots, m$ be a finite collection of such parameters, where $f_j : X \to \mathbb{R}^{n_j}$. Define
    $$
        \bar{F} = \bigcup_{j=1}^m F_j, \quad \bar{\epsilon} = \min_{1 \le j \le m} \epsilon_j,
    $$
    and let $\bar{f} = (f_1, \dots, f_m) : X \to \mathbb{R}^{n_1 + \dots + n_m}$. It is clear that
    $$
        A_{\bar{f}, \bar{\epsilon}, \bar{F}} \subseteq \bigcap_{j=1}^m A_{f_j, \epsilon_j, F_j}.
    $$
    By hypothesis (ii), the set $A_{\bar{f}, \bar{\epsilon}, \bar{F}}$ is non-empty. Therefore, any finite intersection of these sets is non-empty, which proves the claim.

    Returning to the main proof, fix a point $x \in \bigcap_{f, \epsilon, F} A_{f, \epsilon, F}$. We assert that $x$ is a fixed point. Suppose, for the sake of contradiction, that there exists some $g \in G$ such that $g \cdot x \neq x$. Because $X$ is a compact Hausdorff space, Urysohn's Lemma guarantees the existence of a continuous function $f: X \to \mathbb{R}$ such that $f(x) = 0$ and $f(g \cdot x) = 1$. However, this would mean $x \notin A_{f, 1/2, \{g\}}$, which contradicts the fact that $x$ belongs to the intersection of all such sets. Thus, $g \cdot x = x$ for all $g \in G$, and $x$ is indeed a fixed point.
\end{proof}


We use this to prove the following preliminary characterization.



\begin{proposition}
    \label{prop:coset_coloring_equiv}
    Let $S_\infty$ be the Polish group of permutations of $\mathbb{N}$ equipped with the topology of pointwise convergence. If $G \le S_\infty$ is a closed subgroup, then the following are equivalent:
    \begin{enumerate}
        \item[(i)] $G$ is extremely amenable.
        \item[(ii)] For any open subgroup $V$ of $G$, every finite coloring $c: G/V \to \{1, \dots, k\}$ of the set of left-cosets $hV$, and every finite subset $A \subseteq G/V$, there exist $g \in G$ and an index $1 \le i \le k$ such that $c(g \cdot a) = i$ for all $a \in A$. Here, $G$ acts on $G/V$ via the standard left translation $g \cdot hV = ghV$.
    \end{enumerate}
\end{proposition}

\begin{proof}
    (i) $\implies$ (ii). Fix $V, k$, and $c$ as in (ii), and consider the shift action of $G$ on the product space $Y = \{1, \dots, k\}^{G/V}$, given by $g \cdot p(x) = p(g^{-1} \cdot x)$ for $p \in Y$ and $x \in G/V$. This defines a $G$-flow, and $c \in Y$.

    Let $X = \overline{G \cdot c}$ be the orbit closure. By (i), $G$ is extremely amenable, so there exists a fixed point $\gamma \in X$. Since $G$ acts transitively on $G/V$, it is clear that $\gamma: G/V \to \{1, \dots, k\}$ must be a constant function; say $\gamma(a) = i$ for all $a \in G/V$.
    Now, fix a finite subset $A \subseteq G/V$. Because $\gamma \in \overline{G \cdot c}$, there exists some $g \in G$ such that $g^{-1} \cdot c|_A = \gamma|_A$. Consequently, $c(g \cdot a) = \gamma(a) = i$ for all $a \in A$.

    (ii) $\implies$ (i). Observe that (ii) is equivalent to the corresponding statement about the space of right-cosets $Vh$ of $V$, on which $G$ acts by right translation $g \cdot Vh = Vhg^{-1}$; we will utilize this right-coset formulation below. Using Lemma \ref{lem:fixed_point_equiv}, it suffices to show that if $X$ is a $G$-flow, $f: X \to \mathbb{R}^n$ is continuous, $\epsilon > 0$, and $F \subseteq G$ is finite, then there exists $x \in X$ such that $\|f(x) - f(h \cdot x)\| \le \epsilon$ for all $h \in F$.

    By the uniform continuity of the group action at the identity $1_G$, there exists an open neighborhood $V$ of $1_G$ such that for all $h \in V$ and all $x \in X$, we have $\|f(x) - f(h \cdot x)\| \le \epsilon / 3$. Furthermore, since $G$ is a closed subgroup of $S_\infty$, we may assume without loss of generality that $V$ is an open subgroup of $G$ (see, e.g., \cite{BK}, 1.5).

    Next, partition the compact image $f(X) \subseteq \mathbb{R}^n$ into finitely many sets $A_1, \dots, A_k$, each having diameter at most $\epsilon / 3$. Fix an arbitrary point $x_0 \in X$ and define
    $$
        U_i = \{g \in G : f(g \cdot x_0) \in A_i\}.
    $$
    Let $V_i = V U_i$. Since $V_i$ is a union of right-cosets of $V$, it can naturally be viewed as a subset of the right-coset space. Because $\bigcup_{i=1}^k V_i$ covers the entire space, we can define a coloring $c$ of the right-cosets such that $c^{-1}(\{i\}) \subseteq V_i$.

    By hypothesis (ii) (adapted for right-cosets), there exist an index $1 \le i \le k$ and an element $g \in G$ such that $(F \cup \{1_G\})g \subseteq V_i = V U_i$. We claim that setting $x = g \cdot x_0$ satisfies our requirements.

    Indeed, fix $h \in F \cup \{1_G\}$. By construction, $hg \in V U_i$, meaning there exists $v \in V$ such that $v^{-1} h g \in U_i$. By the definition of $U_i$, we have $f(v^{-1} h g \cdot x_0) \in A_i$, which rewrites to $f(v^{-1} h \cdot x) \in A_i$.
    Because $v^{-1} \in V$, the choice of $V$ guarantees that $\|f(v^{-1} h \cdot x) - f(h \cdot x)\| \le \epsilon / 3$. Thus, $f(h \cdot x)$ lies within an $\epsilon / 3$-neighborhood of the set $A_i$.
    Since $1_G \in F \cup \{1_G\}$, the point $f(x)$ also lies within an $\epsilon / 3$-neighborhood of $A_i$. Because the diameter of $A_i$ is at most $\epsilon / 3$, the triangle inequality yields
    $$
        \|f(x) - f(h \cdot x)\| \le \frac{\epsilon}{3} + \frac{\epsilon}{3} + \frac{\epsilon}{3} = \epsilon.
    $$
    This completes the proof.
\end{proof}



Clearly, in Proposition \ref{prop:coset_coloring_equiv}, we can restrict $V$ to any local basis at $1_G$ consisting of open subgroups. In particular, if for each non-empty finite set $F \subseteq \mathbb{N}$ we let
$$
    G_{(F)} = \{g \in G : \forall i \in F(g(i) = i)\}
$$
be the \textit{pointwise stabilizer} of $F$, then the collection $\{G_{(F)} : \emptyset \neq F \subseteq \mathbb{N}, F \text{ is finite}\}$ forms a local basis at $1_G$. Thus, we can restrict $V$ in Proposition \ref{prop:coset_coloring_equiv} to be of the form $G_{(F)}$. Moreover, it suffices to consider only those sets $F$ that belong to a cofinal family (under inclusion) of finite subsets of $\mathbb{N}$.

\paragraph{Remark.}
The proof of Proposition \ref{prop:coset_coloring_equiv} shows that to test for extreme amenability of a closed subgroup $G \le S_\infty$, it is sufficient to find fixed points in compact invariant subsets of the $G$-flow $\{1, \dots, k\}^{G/V}$ for each open subgroup $V \le G$.

\section{Types, Stabilizers, and the Ramsey Property}
\label{sec:ramsey_property}

\paragraph{Further Notation and Terminology.}
For the next result, we introduce some additional concepts.

\subparagraph{Group Actions on Subsets.}
Each subgroup $G \le S_\infty$ acts on the finite subsets of $\mathbb{N}$ in the natural way:
$$
    g \cdot F = \{g(i) : i \in F\}.
$$
For each non-empty finite set $F \subseteq \mathbb{N}$, we define the \textit{setwise stabilizer} of $F$ as
$$
    G_F = \{g \in G : g \cdot F = F\}.
$$
It is clear that the pointwise stabilizer is a subgroup of the setwise stabilizer, $G_{(F)} \le G_F$, and that the index $[G_F : G_{(F)}]$ is finite.

\subparagraph{G-Types.}
The $G$-\textit{type} of a non-empty finite set $F \subseteq \mathbb{N}$ is its orbit under this action, denoted $G \cdot F$. A set $\sigma$ is a $G$-\textit{type} if $\sigma = G \cdot F$ for some such $F$. For two $G$-types $\rho$ and $\sigma$, we define a partial order $\rho \le \sigma$ by the following equivalent conditions:
\begin{itemize}
    \item[$\Leftrightarrow$] $\exists F \in \sigma, \exists F' \in \rho$ such that $F' \subseteq F$.
    \item[$\Leftrightarrow$] $\forall F \in \sigma, \exists F' \in \rho$ such that $F' \subseteq F$.
    \item[$\Leftrightarrow$] $\forall F' \in \rho, \exists F \in \sigma$ such that $F' \subseteq F$.
\end{itemize}

\subparagraph{The Logic Action.}
Given a signature $L = \{R_i\}_{i \in I} \cup \{f_j\}_{j \in J}$, we denote by $X_L$ the space of all $L$-structures with universe $\mathbb{N}$. Thus,
$$
    X_L = \prod_{i \in I} 2^{(\mathbb{N}^{n_i})} \times \prod_{j \in J} \mathbb{N}^{(\mathbb{N}^{m_j})},
$$
where $n_i$ and $m_j$ are the arities of the symbols. If the signature $L$ is purely relational (i.e., $J = \emptyset$), then $X_L$ is a compact space, homeomorphic to $2^\mathbb{N}$.

The group $S_\infty$ acts canonically on $X_L$. Given a structure $\mathbf{a} = \langle \mathbb{N}, \{R^\mathbf{a}_i\}, \{f^\mathbf{a}_j\} \rangle$, we define its image under $g \in S_\infty$ as $g \cdot \mathbf{a} = \mathbf{b} = \langle \mathbb{N}, \{R^\mathbf{b}_i\}, \{f^\mathbf{b}_j\} \rangle$, where
$$
    R^\mathbf{b}_i(a_1, \dots, a_{n_i}) \iff R^\mathbf{a}_i(g^{-1}(a_1), \dots, g^{-1}(a_{n_i})),
$$
$$
    f^\mathbf{b}_j(a_1, \dots, a_{m_j}) = g(f^\mathbf{a}_j(g^{-1}(a_1), \dots, g^{-1}(a_{m_j}))).
$$
This action, called the \textit{logic action}, is continuous, and $g$ serves as an isomorphism from $\mathbf{a}$ to $\mathbf{b}$. In particular, if $L$ is relational, the action of any $G \le S_\infty$ on $X_L$ is a $G$-flow.

As a key example, consider the language $L = \{<\}$, where $<$ is a binary relation symbol. Let $\mathrm{LO} \subseteq X_L$ be the compact, $S_\infty$-invariant subset consisting of all linear orderings on $\mathbb{N}$. For any $G \le S_\infty$, $\mathrm{LO}$ is a subflow of the $G$-flow $X_L$. We say that $G$ \textit{preserves an ordering} if this subflow has a fixed point; that is, if there exists a linear order $\prec$ on $\mathbb{N}$ such that for every $g \in G$, $a \prec b \Leftrightarrow g(a) \prec g(b)$.


We now have:


\begin{proposition}
    \label{prop:stabilizer_ramsey_equiv}
    Let $G \le S_\infty$ be a closed subgroup. Then the following are equivalent:
    \begin{enumerate}
        \item[(i)] $G$ is extremely amenable.
        \item[(ii)]
              \begin{enumerate}
                  \item[\textbf{(a)}] For every non-empty finite set $F \subseteq \mathbb{N}$, the pointwise and setwise stabilizers coincide: $G_{(F)} = G_F$; and
                  \item[\textbf{(b)}] For any two $G$-types $\rho, \sigma$ with $\rho \le \sigma$, and every finite coloring $c: \rho \to \{1, \dots, k\}$, there exist an index $1 \le i \le k$ and a set $F \in \sigma$ such that $c(F') = i$ for all $F' \subseteq F$ with $F' \in \rho$.
              \end{enumerate}
        \item[(iii)]
              \begin{enumerate}
                  \item[\textbf{(a)'}] $G$ preserves an ordering; and
                  \item[\textbf{(b)}] The same condition as in (ii)(b).
              \end{enumerate}
    \end{enumerate}
\end{proposition}

\begin{proof}
    (i) $\implies$ (iii). We first consider (a)'. Since $G$ is extremely amenable and the space of linear orderings $\mathrm{LO}$ is a $G$-flow, there must exist a fixed point. This fixed point precisely means that $G$ preserves an ordering.

    Next, we prove (b). Fix $G$-types $\rho \le \sigma$ and a coloring $c: \rho \to \{1, \dots, k\}$. Suppose $\rho = G \cdot F'$, and let $V = G_{(F')}$. By (a)' (which implies (a)), we have $V = G_{(F')} = G_{F'}$. Because the setwise and pointwise stabilizers are equal, we can identify the coset space $G/V$ with the orbit $G \cdot F' = \rho$.
    Now, apply condition (ii) of Proposition \ref{prop:coset_coloring_equiv} to the subgroup $V$, the coloring $c$, and the finite subset $A = \{F'_0 \subseteq F_0 : F'_0 \in \rho\}$, where $F_0$ is a fixed representative in $\sigma$. This yields an index $1 \le i \le k$ and an element $g \in G$ such that $c(g \cdot F'_0) = i$ for all $F'_0 \in A$.
    Let $F = g \cdot F_0 \in \sigma$. If $F' \subseteq F$ is such that $F' \in \rho$, then $g^{-1} \cdot F' = F'_0 \subseteq F_0$, meaning $F'_0 \in A$. Thus, $c(F') = c(g \cdot F'_0) = i$, which completes the proof of (b).

    (iii) $\implies$ (ii). It is clear that (a)' $\implies$ (a), because if a permutation group preserves a linear ordering, any element that fixes a finite set setwise must necessarily fix its elements pointwise. Condition (b) is identical in both.

    (ii) $\implies$ (i). We will verify condition (ii) of Proposition \ref{prop:coset_coloring_equiv} for an open subgroup $V$ of the form $V = G_{(F)} = G_F$, where $F \subseteq \mathbb{N}$ is a non-empty finite set. If $V = G_F$, then the space $G/V$ can be naturally identified with the $G$-type $\rho = G \cdot F$.

    Fix a finite coloring $c: \rho \to \{1, \dots, k\}$ and a finite subset $A \subseteq \rho$. Define $F_0 = \bigcup A$ and let $\sigma = G \cdot F_0$. It is clear from the construction that $\rho \le \sigma$.
    By condition (ii)(b), there exist an index $1 \le i \le k$ and an element $g \in G$ such that for all $F' \subseteq g \cdot F_0$ with $F' \in \rho$, we have $c(F') = i$.

    Since every $F \in A$ satisfies $F \in \rho$ and $F \subseteq F_0$, it follows that $g \cdot F \in \rho$ and $g \cdot F \subseteq g \cdot F_0$. Therefore, substituting $F' = g \cdot F$ into the above property yields $c(g \cdot F) = i$ for all $F \in A$. This satisfies the sufficient condition for extreme amenability given in Proposition \ref{prop:coset_coloring_equiv}.
\end{proof}



We will now use a compactness argument to put this characterization in its final form. It will be convenient to first introduce the following notation.

\begin{definition}
    \label{def:ramsey_property}
    Let $G \le S_\infty$ be a closed subgroup, and let $\rho \le \sigma$ be $G$-types. For any set $F \in \sigma$, we define the collection of its subsets of type $\rho$ as:
    $$
        \binom{F}{\rho} = \{F' \subseteq F : F' \in \rho\}.
    $$

    Let $\rho \le \sigma \le \tau$ be $G$-types, and let $k \ge 2$ be an integer. We write the partition relation
    $$
        \tau \to (\sigma)^\rho_k
    $$
    if for every $F \in \tau$ and every finite coloring $c: \binom{F}{\rho} \to \{1, \dots, k\}$, there exists a subset $F_0 \in \binom{F}{\sigma}$ that is \textit{homogeneous}. That is, $c$ is monochromatic on $\binom{F_0}{\rho}$: there exists an index $1 \le i \le k$ such that $c(F') = i$ for all $F' \in \binom{F_0}{\rho}$.

    \textit{(Note: Because $\tau$ is a single $G$-orbit, asserting this condition holds for \textbf{every} $F \in \tau$ is equivalent to asserting it holds for \textbf{some} $F \in \tau$.)}

    We say that the group $G$ has the \textbf{Ramsey property} if for all $G$-types $\rho \le \sigma$ and every integer $k \ge 2$, there exists a $G$-type $\tau \ge \sigma$ such that
    $$
        \tau \to (\sigma)^\rho_k.
    $$
\end{definition}

We now have

\begin{theorem}
    \label{thm:ea_ordering_ramsey_equiv}
    Let $G \le S_\infty$ be a closed subgroup. Then the following are equivalent:
    \begin{enumerate}
        \item[(i)] $G$ is extremely amenable.
        \item[(ii)]
              \begin{enumerate}
                  \item[\textbf{(a)}] $G$ preserves an ordering; and
                  \item[\textbf{(b)}] $G$ has the Ramsey property.
              \end{enumerate}
    \end{enumerate}
\end{theorem}

\begin{proof}
    (i) $\implies$ (ii). The fact that (i) implies (a) was already established in Proposition \ref{prop:stabilizer_ramsey_equiv}. To prove (b), we first note that by a standard induction argument, it is sufficient to verify the Ramsey property defined in Definition \ref{def:ramsey_property} for the case $k=2$.

    Suppose, towards a contradiction, that there exist $G$-types $\rho \le \sigma$ for which there is no $\tau \ge \sigma$ satisfying $\tau \to (\sigma)^\rho_2$. Fix a representative set $F_0 \in \sigma$. By our assumption, for every finite set $E \supseteq F_0$, there exists a coloring
    $$
        c_E: \binom{E}{\rho} \to \{1, 2\}
    $$
    that admits no homogeneous set $F \in \binom{E}{\sigma}$.

    Let $I$ be the set of all non-empty finite subsets of $\mathbb{N}$. The collection of sets of the form $\{E \in I : F \subseteq E\}$ for finite $F \subseteq \mathbb{N}$ clearly has the finite intersection property. Thus, we can choose an ultrafilter $U$ on $I$ that extends this collection.

    For each $D \in \rho$, exactly one of the sets $\{E \in I : E \supseteq D \cup F_0 \text{ and } c_E(D) = 1\}$ or $\{E \in I : E \supseteq D \cup F_0 \text{ and } c_E(D) = 2\}$ belongs to $U$. We can therefore define a limit coloring $c: \rho \to \{1, 2\}$ by setting $c(D) = i$ if and only if
    $$
        \{E \in I : E \supseteq D \cup F_0 \text{ and } c_E(D) = i\} \in U.
    $$

    By Proposition \ref{prop:stabilizer_ramsey_equiv}(ii)(b), there exists a set $F \in \sigma$ such that $c$ is monochromatic on $\binom{F}{\rho}$, say with constant value $i \in \{1, 2\}$. For each $D \in \binom{F}{\rho}$, define the set
    $$
        A_D = \{E \in I : E \supseteq F \cup F_0 \text{ and } c_E(D) = c(D) = i\}.
    $$
    By the definition of our limit coloring, $A_D \in U$ for all $D \in \binom{F}{\rho}$. Because $\binom{F}{\rho}$ is a finite collection and $U$ is closed under finite intersections, the intersection $\bigcap_{D \in \binom{F}{\rho}} A_D$ also belongs to $U$ and is therefore non-empty.

    Pick an element $E$ in this intersection. Then $E \supseteq F_0$, meaning the coloring $c_E$ is defined, and $E \supseteq F$. Furthermore, for every $D \in \binom{F}{\rho}$, we have $c_E(D) = i$. This implies that $F \in \binom{E}{\sigma}$ is a homogeneous set for $c_E$, which contradicts the original choice of the coloring $c_E$.

    (ii) $\implies$ (i). By Proposition \ref{prop:stabilizer_ramsey_equiv}(iii), it is sufficient to verify condition (b) of Proposition \ref{prop:stabilizer_ramsey_equiv}(ii). Let $\rho \le \sigma$ be $G$-types, and let $c: \rho \to \{1, \dots, k\}$ be a finite coloring. Because $G$ has the Ramsey property, there exists a $G$-type $\tau \ge \sigma$ such that $\tau \to (\sigma)^\rho_k$.

    Choose any $E \in \tau$. The restriction of $c$ to $\binom{E}{\rho}$ is a valid $k$-coloring, so by the partition relation, there exists a homogeneous set $F \in \binom{E}{\sigma}$. This means $F \in \sigma$, $F \subseteq E$, and $c$ is constant on all $F' \in \binom{F}{\rho}$, which trivially satisfies the required condition.
\end{proof}


\begin{remark}
    \label{rem:cofinal_ramsey}
    Let $G \le S_\infty$. We call a set $T$ of $G$-types \textit{cofinal} if for every $G$-type $\rho$, there exists a type $\sigma \in T$ such that $\rho \le \sigma$. It is not hard to see that Theorem \ref{thm:ea_ordering_ramsey_equiv} still holds if, in Definition \ref{def:ramsey_property} (the Ramsey property for $G$), we restrict the $G$-types to belong to any given cofinal set of $G$-types.
\end{remark}

\section{Characterization via Fraïssé Limits}
\label{sec:fraisse_characterization}



We will finally tie up the extreme amenability of automorphism groups with the structural Ramsey theory discussed in Section 3.

Let $L$ be a signature with a distinguished binary relation symbol $<$ (and perhaps other symbols). An \textit{order structure} for $L$ is an $L$-structure $\mathbf{A}$ in which $<^\mathbf{A}$ is a linear ordering. If $\mathcal{K}$ is a class of $L$-structures, we say that $\mathcal{K}$ is an \textit{order class} if every $\mathbf{A} \in \mathcal{K}$ is an order structure.

Recall that, up to topological group isomorphism, the closed subgroups of $S_\infty$ are exactly the same as the automorphism groups of countable structures. Equivalently, they are exactly the Polish groups that admit a countable neighborhood basis at the identity consisting of open subgroups (see \cite[Section 1.5]{BK}). The next result provides a characterization of the extremely amenable groups within this class.



\begin{theorem}
    \label{thm:main_characterization}
    Let $G \le S_\infty$ be a closed subgroup. Then the following are equivalent:
    \begin{enumerate}
        \item[(i)] $G$ is extremely amenable.
        \item[(ii)] $G = \operatorname{Aut}(\mathbf{A})$, where $\mathbf{A}$ is the Fraïssé limit of a Fraïssé order class with the Ramsey property.
    \end{enumerate}
\end{theorem}

\begin{proof}
    (i) $\implies$ (ii): Let $\mathbf{A}_G = \langle \mathbb{N}, \dots \rangle$ be the canonical relational structure for $G$ induced by the logic action. As noted in Section 2, $\mathbf{A}_G$ is ultrahomogeneous. Since $G$ is extremely amenable, it preserves a linear order $\prec$ on $\mathbb{N}$. Let $L$ be the signature of $\mathbf{A}_G$ augmented with a new binary relation symbol $<$. Let $\mathbf{A}$ be the expansion of $\mathbf{A}_G$ in which $<^\mathbf{A} = \prec$.

    The automorphism group of this expanded structure is precisely $G$, since any automorphism must preserve all relations, including the new ordering. Thus, $\operatorname{Aut}(\mathbf{A}) = G$, which implies that $\mathbf{A}$ is also ultrahomogeneous. As the signature of $\mathbf{A}_G$ is relational, the structure $\mathbf{A}$ is locally finite. It follows that its age, $\mathcal{K} = \operatorname{Age}(\mathbf{A})$, is a Fraïssé class. Since every structure in $\mathcal{K}$ is a finite substructure of the order structure $\mathbf{A}$, $\mathcal{K}$ is a Fraïssé order class.

    Now, we connect the Ramsey property of $G$ with that of $\mathcal{K}$. By the ultrahomogeneity of $\mathbf{A}$, a $G$-type of a finite set corresponds precisely to the isomorphism type of the substructure of $\mathbf{A}$ on that set. Under this correspondence, the Ramsey property for $G$ (as defined in Definition \ref{def:ramsey_property}) is equivalent to the structural Ramsey property for $\operatorname{Age}(\mathbf{A})$. Since $G$ is extremely amenable, Theorem \ref{thm:ea_ordering_ramsey_equiv} implies that $G$ has the Ramsey property, and therefore $\operatorname{Age}(\mathbf{A})$ has the Ramsey property. This completes the proof of this direction.

    \medskip

    (ii) $\implies$ (i): Assume $G = \operatorname{Aut}(\mathbf{A})$, where $\mathbf{A}$ is the Fraïssé limit of a Fraïssé order class $\mathcal{K}$ that has the Ramsey property.

    First, since $\mathbf{A}$ is an order structure by hypothesis, its automorphism group $G$ must preserve the linear ordering on $\mathbf{A}$. This satisfies condition (a) of Theorem \ref{thm:ea_ordering_ramsey_equiv}(ii).

    Second, because $\mathbf{A}$ is a locally finite structure, the $G$-types corresponding to the isomorphism types of finite substructures in $\operatorname{Age}(\mathbf{A})$ form a cofinal set of all $G$-types. The assumption that $\mathcal{K} = \operatorname{Age}(\mathbf{A})$ has the Ramsey property translates directly to the fact that $G$ satisfies the Ramsey partition relation $\tau \to (\sigma)^\rho_k$ for this cofinal family of types. By Remark \ref{rem:cofinal_ramsey}, this is sufficient to conclude that $G$ has the Ramsey property, which satisfies condition (b) of Theorem \ref{thm:ea_ordering_ramsey_equiv}(ii).

    Since both conditions of Theorem \ref{thm:ea_ordering_ramsey_equiv}(ii) are met, we conclude that $G$ is extremely amenable.
\end{proof}



We make explicit the following fact observed in the preceding proof.

\begin{theorem}
    \label{thm:fraisse_ramsey_equiv}
    Let $\mathcal{K}$ be a Fraïssé order class and let $\mathbf{A} = \operatorname{Flim}(\mathcal{K})$ be its Fraïssé limit. Then the following are equivalent:
    \begin{enumerate}
        \item[(i)] $\operatorname{Aut}(\mathbf{A})$ is extremely amenable.
        \item[(ii)] $\mathcal{K}$ has the Ramsey property.
    \end{enumerate}
\end{theorem}

\begin{proof}
    The equivalence was established during the proof of Theorem \ref{thm:main_characterization}.
\end{proof}
